\documentclass[10pt,a4paper]{amsart}
\usepackage[margin=19mm]{geometry}
\usepackage{amsmath,amssymb,mathtools}
\usepackage[scr=rsfs,cal=euler]{mathalfa}
\usepackage{xfrac}
\usepackage[unicode,hidelinks,bookmarks=false]{hyperref}
\newcommand{\ie}{i.e.\ }
\newcommand{\cf}{cf.\ }
\newcommand{\resp}{respectively\ }
\newcommand{\Subsection}{\subsection}
\providecommand{\nobreakdash}{\nobreak\hbox{-}}
\setlength{\emergencystretch}{2em}
\multlinegap0pt
\usepackage{bm}
\usepackage{bbm}
\usepackage{dsfont}
\usepackage{xspace}
\usepackage{cancel}
\usepackage{mathtools}
\usepackage{amsthm}
\makeatletter
\@ifundefined{theorem}{\newtheorem{theorem}{Theorem}}{}
\@ifundefined{lemma}{\newtheorem{lemma}[theorem]{Lemma}}{}
\makeatother
\newcommand{\R}{\mathbb R}
\newcommand{\Z}{\mathbb Z}
\title[Is the space of reachable particle configurations dense?]{Is the space of reachable particle configurations dense?}
\author{Janos Pach, Gabor Tardos}
\date{Source: arXiv:2507.22471v1 (CC BY 4.0)}
\begin{document}
\maketitle

\noindent\emph{Source and licence.} Is the space of reachable particle configurations dense?. Version arXiv:2507.22471v1 (CC BY 4.0), distributed under the Creative Commons Attribution 4.0 International licence, as recorded in the arXiv licence metadata for this exact version and shown on the abstract page https://arxiv.org/abs/2507.22471v1. Full rights notice: licenses/arxiv-2507.22471-CC-BY-4.0.txt.\par\smallskip

\noindent\textit{Editorial scope.} Counted unit: the Lemma whose source label is \texttt{None} in the article (source0.tex lines 156--170, source label \texttt{None}), delivered together with the article's own complete original proof of it (source0.tex lines 172--194, one \texttt{proof} environment of 23 lines). No mathematics is written, repaired or re-derived: statement and proof are the article's own text line for line. Required context retained uncounted: none. Every definition, display and label of those lines is retained unchanged. Omitted from this package and not redistributed: 1--155, 171--171, 195--437. Nothing inside a retained excerpt is omitted or altered: every retained body is delivered exactly as the article's lines are written, so no omission receipt of any kind accompanies this package. Citations: the retained excerpts carry no citation command at all, so no bibliography entry of the article is transcribed and no reference is redistributed. Reference adaptation: none was needed and none was made; every reference inside the delivered unit resolves inside it, so no unresolved reference or citation is produced. Printed slips kept literal: no printed word, symbol, hypothesis or number of the counted statement or of its proof was altered. Layout and class: the article's own class is \texttt{article} with packages including amsthm, amsmath, amssymb, enumitem; the wrapper is the lane's pinned \texttt{amsart} editorial wrapper at 10pt on A4, which restates the theorem-like environments the retained text needs and adds the article's own macro definitions that the retained text needs (\texttt{\textbackslash R}, \texttt{\textbackslash Z}), with extra wrapper packages (bm, bbm, dsfont, xspace, cancel, mathtools). The delivered numbering is the unit's own. Assets: no retained span contains a figure environment, a graphics-inclusion command, a PDF-page inclusion, a table or a TikZ picture; no image file is copied into this package, the package declares no asset dependency and no image file is redistributed with it. Licence: version arXiv:2507.22471v1 of this article is distributed under the Creative Commons Attribution 4.0 International licence, as recorded in the arXiv licence metadata for this exact version and shown on the abstract page https://arxiv.org/abs/2507.22471v1. A full rights notice is delivered at licenses/arxiv-2507.22471-CC-BY-4.0.txt. Characters: every retained excerpt is pure ASCII, so the delivered engine, XeLaTeX with its default Latin Modern text font, reports no missing character.
\begin{lemma}\label1
\begin{itemize}
\item[(i)]
Let $n$ be a positive integer and $v\in\R^n$. There exists a $n\times n$ step matrix $A$ such that the vector $Av$ is shorter than $v$, unless every entry of $v$ is equal to $0$, $a$, or $-a$, for some $a\in\R$.
\item[(ii)]
Let $n$ be a positive integer and $v\in\Z^n$ a nonzero vector. There exists a good $n\times n$ matrix $A$ such that the only entries of $Av$ are $a$, $-a$, and $0$, where $a$ is the greatest common divisor of the entries of $v$.
\item[(iii)]
A square matrix $A$ is good if and only if
\begin{enumerate}
\item $\det(A)=1$;
\item the non-diagonal entries of $A$ are even integers; and
\item the diagonal entries of $A$ are integers congruent to $1$ modulo $4$.
\end{enumerate}
\end{itemize}
\end{lemma}

\begin{proof}
To see part~(i), consider a vector $v\in\R^n$ with two entries $a$ and $b$ satisfying $|b|>|a|>0$. If we multiply $v$ by a suitable step matrix $A$ we can change its entry $b$ to either $b+2a$ or $b-2a$, without altering any other entries. One of these changes decreases the Euclidean norm of the vector, that is, we have $||Av||<||v||.$
\medskip

For part~(ii), we take a vector $v\in\Z^n$ and apply part~(i) to it, repeatedly. The length of the resulting integer vector decreases in every step. Hence, we must get stuck after finitely many steps at an integer vector $Av$, where $A$ is good and all nonzero entries of $Av$ have the same absolute value. During this process, the greatest common divisor, $a$, of the entries of the vector does not change, therefore all entries of $Av$ must be $a$, $-a$ or $0$, as claimed.
\medskip

Notice that every step matrix satisfies all three conditions listed in part~(iii). Moreover, it is easy to see that the set of matrices $A$ satisfying these conditions form a subgroup. This proves the ``only if'' direction of part~(iii).
\medskip

It remains to prove the ``if'' direction of part~(iii). We proceed by induction on $n$. The case $n=1$ is trivial (only the identity matrix satisfies the conditions), so we assume $n>1$ and also that the statement holds for smaller matrices.

\smallskip

Consider an $n\times n$ integer matrix $X$ satisfying the three conditions in the statement. We will multiply $X$ from the left with good matrices. Note that the matrices we obtain during this process must all satisfy our three conditions.
Our goal is to transform the matrix into simpler and simpler matrices. When we arrive at the identity matrix, we have succeeded in proving that $X$ is good.

First, we apply part~(ii) of the lemma to the last column $v$ of $X$. Note that the greatest common divisor of the entries of $v$ divides $\det(X)=1$. Thus, we can find a good matrix $A$ with all entries of $Av$ being $1$, $-1$ or $0$. Note that the last column of $AX$ is $Av$, and $AX$ satisfies the divisibility conditions, so its last column, $Av$, must be equal to $e_n$, the last standard basis vector of $\R^n$ (i.e., all-zero with a single $1$ entry in the last place).

Let $B$ denote the submatrix of $AX$ formed by the first $n-1$ rows and columns. It must satisfy the same divisibility conditions as $AX$ does. Expanding the determinant of $AX$ along its last column, we obtain that $\det(B)=\det(AX)=1.$ By the induction hypothesis, $B$ must be good. So, in particular, $B^{-1}$ can be written as product of step matrices: $B^{-1}=\prod_{i=1}^kC_i$. Let $f(C)$ stand for the extension of the $(n-1)\times(n-1)$ matrix $C$ by a last row and column containing only zeros except for the $1$ in the diagonal. The function $f$ preserves multiplication, so we have $f(B^{-1})=\prod_{i=1}^kf(C_i)$. Here, $f(C_i)$ is a step matrix for every $i$, so $f(B^{-1})$ is also a good matrix.

Let us consider now the matrix $f(B^{-1})AX$. It differs from the identity matrix only in the first $n-1$ entries in the last row. These entries must be even, as our divisibility conditions must be satisfied. The $i$\/th entry in the last row can then be get rid of by multiplying it with the appropriate power of the step matrix having $2$ in the same position. This completes our project to obtain the identity matrix by multiplying $X$ with good matrices and, therefore, shows that $X$ is good.
\end{proof}

\end{document}
